%%=============================================================================
%% Samenvatting
%%=============================================================================

% TODO: De "abstract" of samenvatting is een kernachtige (~ 1 blz. voor een
% thesis) synthese van het document.
%
% Een goede abstract biedt een kernachtig antwoord op volgende vragen:
%
% 1. Waarover gaat de bachelorproef?
% 2. Waarom heb je er over geschreven?
% 3. Hoe heb je het onderzoek uitgevoerd?
% 4. Wat waren de resultaten? Wat blijkt uit je onderzoek?
% 5. Wat betekenen je resultaten? Wat is de relevantie voor het werkveld?
%
% Daarom bestaat een abstract uit volgende componenten:
%
% - inleiding + kaderen thema
% - probleemstelling
% - (centrale) onderzoeksvraag
% - onderzoeksdoelstelling
% - methodologie
% - resultaten (beperk tot de belangrijkste, relevant voor de onderzoeksvraag)
% - conclusies, aanbevelingen, beperkingen
%
% LET OP! Een samenvatting is GEEN voorwoord!

%%---------- Nederlandse samenvatting -----------------------------------------
%
% TODO: Als je je bachelorproef in het Engels schrijft, moet je eerst een
% Nederlandse samenvatting invoegen. Haal daarvoor onderstaande code uit
% commentaar.
% Wie zijn bachelorproef in het Nederlands schrijft, kan dit negeren, de inhoud
% wordt niet in het document ingevoegd.

\IfLanguageName{english}{%
\selectlanguage{dutch}
\chapter*{Samenvatting}
\lipsum[1-4]
\selectlanguage{english}
}{}

%%---------- Samenvatting -----------------------------------------------------
% De samenvatting in de hoofdtaal van het document

\chapter*{\IfLanguageName{dutch}{Samenvatting}{Abstract}}

Betrouwbare toegang tot digitale bestanden veronderstelt doorgaans een stabiele
internetverbinding -- een aanname die niet standhoudt op zee, in afgelegen gebieden of
tijdens rampen waarbij de netwerkinfrastructuur wegvalt. Commerciële alternatieven zoals
satellietverbindingen zijn voor particulieren en kleine organisaties vaak financieel of
praktisch ontoereikend. Radioverbindingen worden in de literatuur herhaaldelijk naar
voren geschoven als infrastructuuronafhankelijk alternatief, maar hun specifieke
toepasbaarheid voor toegang tot een \emph{eigen, persoonlijke fileserver} -- in
tegenstelling tot de courant onderzochte toepassingen zoals noodcommunicatie of
professionele backhaul -- was tot dusver niet grondig onderzocht.
 
Deze bachelorproef onderzoekt daarom hoe een verbinding met een private fileserver via
een radioverbinding kan worden opgezet, en wat de technische randvoorwaarden en
maximale afstandsbeperkingen zijn voor een werkende implementatie. De doelgroep bestaat
uit maritieme operatoren, hulpverleningsorganisaties, onderzoekers in afgelegen gebieden
en radioamateurs. Het onderzoek streeft een werkende proof-of-concept na, beoordeeld aan
de hand van meetbare criteria: doorvoersnelheid, packet loss, verbindingsstabiliteit en
operationele afstand.
 
Op basis van een literatuurstudie naar signaalpropagatie, modulatietechnieken en
bestaande radioprotocollen (Winlink, AX.25) is een architectuur geselecteerd die SFTP
combineert met een 5.6~GHz point-to-point-straalverbinding binnen de licentievrije
ISM-band. Omdat de voorziene fysieke apparatuur niet binnen het beschikbare
tijdsbestek kon worden aangekocht, is deze architectuur gerealiseerd als een volledig
geautomatiseerde, virtuele netwerkemulatie: een dumbbell-topologie van drie virtuele
machines, waarbij een centrale link-emulator via \texttt{tc/netem} het gedrag van de
straalverbinding nabootst. Om de emulatie zo realistisch mogelijk te houden, zijn de
kanaalparameters niet vrij gekozen, maar afgeleid uit officiële productspecificaties,
aangevuld met een Gilbert-Elliot-burstverliesmodel voor willekeurig variërende
weersomstandigheden en een stochastische kans op een wettelijk verplichte
DFS-kanaalonderbreking -- beide onafhankelijk van de afstand. Voor elk van drie
afstandscategorieën (0--15~km, 15--30~km en 30+~km) zijn honderd onafhankelijke
metingen uitgevoerd, in totaal 300 trials, waarna de resultaten statistisch
geanalyseerd en per categorie getoetst zijn aan de vooraf gedefinieerde
netwerkrequirements.
 
De resultaten tonen een duidelijk, oplopend patroon. Binnen het optimale bereik
(0--15~km) voldoet de architectuur ruim aan de norm voor packet loss (gemiddeld
0{,}35\% onder normale omstandigheden), met overschrijdingen die vrijwel uitsluitend het
rechtstreekse gevolg zijn van een verplichte DFS-onderbreking. Vanaf de horizonlimiet
(15--30~km) schiet de courante kanaalkwaliteit zelf reeds structureel tekort (gemiddeld
5{,}27\%), en voorbij het bereik dat als onbetrouwbaar werd ingeschat (30+~km) is die
tekortkoming vrijwel universeel (20{,}38\%). Onafhankelijk van de afstand toont de
architectuur zich wel steeds in staat om een bestandsoverdracht te hervatten na een
DFS-onderbreking, wat de robuustheid van de gekozen combinatie van SFTP en TCP
bevestigt.
 
Deze bachelorproef concludeert dat radiogebaseerde toegang tot een private fileserver
technisch haalbaar is binnen een zorgvuldig afgebakend afstandsbereik van ongeveer
15~km, met een geleidelijke overgang naar onbetrouwbaarheid daarna, in plaats van een
enkele harde grens. Naast dit inhoudelijke antwoord levert het onderzoek een generiek
herbruikbare, statistisch onderbouwde testmethodologie op. De belangrijkste beperking
is dat de resultaten voortkomen uit een netwerkemulatie op basis van plausibele,
literatuurgebaseerde parameters, niet uit fysieke veldmetingen; fysieke validatie van de
architectuur, zodra geschikte apparatuur beschikbaar is, vormt dan ook de voornaamste
aanbeveling voor vervolgonderzoek.
